%%%%%%%%%%  scan idx 302  |  folio 289  %%%%%%%%%%
Supposons maintenant qu'on soit dans la situation (2.11.1) (resp. (2.11.2)), avec \(f\) de tor-dimension finie. Alors (III 4.8.1 et 5.9 b)) le foncteur \(Rf_*\) envoie \(D(X)_{\uparf}\) dans \(D(Y)_{\uparf}\), et par suite définit un homomorphisme
\begin{equation}\tag{2.12.3}
 f_*:K^{\bullet}(X)\longrightarrow K^{\bullet}(Y).
\end{equation}
Celui-ci n'est pas en général un homomorphisme d'anneaux. Cependant, pour \(y\in K_{\bullet}(Y)\) (resp. \(K^{\bullet}(Y)\)) et \(x\in K^{\bullet}(X)\), on a la formule
\begin{equation}\tag{2.12.4}
 f_*(f^*(y)x)=y f_*(x),
\end{equation}
où \(f^*\) est l'homomorphisme (2.12.1) (resp. (2.7.1)). La démonstration est analogue à celle de (2.11.1.2). Si \(g:Y\to Z\) est un morphisme satisfaisant aux mêmes hypothèses que \(f\), les homomorphismes (2.12.3) vérifient la formule de transitivité~:
\begin{equation}\tag{2.12.3.1}
 g_*f_*=(gf)_*.
\end{equation}

L'homomorphisme (2.12.3) est l'un des ingrédients essentiels du théorème de Riemann--Roch, aussi bien en Géométrie Algébrique qu'en Géométrie Analytique (pour des indications dans le cas analytique, voir (Exp. 0, 5)).

Supposons que \(X\) et \(Y\) soient des schémas séparés et divisoriels (II 2.2.5) (ce qui est le cas par exemple si \(X\) et \(Y\) possèdent des faisceaux inversibles amples). Alors (2.9 (i)) \(K^{\bullet}(X)\) (resp. \(K^{\bullet}(Y)\)) s'identifie à \(K^{\bullet}(X)_{\naif}\) (resp. \(K^{\bullet}(Y)_{\naif}\)), et (2.12.3) définit donc un homomorphisme
\begin{equation}\tag{2.12.3'}
 f_*:K^{\bullet}(X)_{\naif}\longrightarrow K^{\bullet}(Y)_{\naif}.
\end{equation}
Comme il a déjà été expliqué dans (Exp. 0, 3.1), la définition de (2.12.3') n'était nullement évidente à priori. Elle constitue même, dans l'optique de notre Séminaire, l'une des motivations principales des
sorites de Exp. I à Exp. IV.

